\documentclass[12pt,a4paper]{article}

\usepackage[UTF8]{ctex}
\usepackage[a4paper,margin=1in]{geometry}
\usepackage{amsmath,amssymb,amsthm,mathtools}
\usepackage{hyperref}
\usepackage{enumitem}
\usepackage{booktabs}

\hypersetup{
    colorlinks=true,
    linkcolor=blue,
    urlcolor=blue,
    citecolor=blue
}

\setlength{\parindent}{2em}
\setlength{\parskip}{0.5em}

\title{Gentry09 中 \(a\) 与 \(a'\) 的关系：\\
从格陪集到模 \(2\) 正确性的精确整理}
\author{}
\date{}

\newtheorem{definition}{定义}
\newtheorem{proposition}{命题}
\newtheorem{remark}{注}
\newtheorem{example}{例}

\begin{document}

\maketitle

\tableofcontents
\newpage

\section{引言}

在学习 Gentry09 时，一个很容易被说糊涂、但实际上又非常关键的问题是：

\begin{quote}
解密过程中得到的量 \(a'\)，和加密时最初构造的
\[
a=m+2e
\]
到底是什么关系？
\end{quote}

很多讲解会直接说“解密恢复出 \(a\)”；但如果严格追问，这句话其实并不总是无条件成立。另一方面，如果只说“\(a'\) 和 \(a\) 模同一个格同余”，这又可能太弱，不足以推出消息比特一定恢复正确。

因此，这个问题的关键不在于泛泛地说“差一个格向量”，而在于：

\begin{enumerate}[label=\arabic*.]
    \item 先精确写出 \(a'\) 与 \(a\) 的代数关系；
    \item 再判断这种关系是否足以推出模 \(2\) 后消息不变；
    \item 最后区分哪些结论是\textbf{必要的}，哪些只是\textbf{更强但非必须的}。
\end{enumerate}

本文只做这一件事：专门把 \(a\) 与 \(a'\) 的关系讲清楚。

\section{记号与问题设置}

\subsection{加密前构造的短偏移}

Gentry09 的简化描述中，消息为
\[
m\in\{0,1\},
\]
并先构造
\[
a = m + 2e,
\]
其中 \(e\) 是小噪声。

这里 \(a\) 的作用不是格点，而是一个承载消息的\textbf{短偏移}。因为 \(2e\) 为偶数项，所以
\[
a \bmod 2 = m.
\]

因此，最终解密只要能恢复一个与 \(a\) 在模 \(2\) 意义下等价的量，就足以恢复消息。

\subsection{格与密钥}

设
\[
L = L(B_{pk}) = L(B_{sk})
\]
表示同一个格，其中

\begin{itemize}[leftmargin=2em]
    \item \(B_{pk}\) 是公钥劣质基；
    \item \(B_{sk}\) 是私钥优质基。
\end{itemize}

加密与解密都围绕这个格展开。

\subsection{密文与解密输出}

可将密文写成
\[
c = \ell + a,
\qquad \ell\in L.
\]
这里 \(\ell\) 表示加密过程中所引入的格点部分。

解密时，私钥算法根据 \(c\) 再恢复一个格点 \(\ell'\in L\)，并输出
\[
a' = c - \ell'.
\]

问题就在于：这个 \(a'\) 与原始 \(a\) 到底是什么关系？

\section{从定义直接推出的关系}

这是最基本、最不会错的一步。

由
\[
c=\ell+a
\]
以及
\[
a' = c-\ell'
\]
立刻得到
\[
a' = (\ell+a)-\ell' = a + (\ell-\ell').
\]

令
\[
\lambda = \ell-\ell',
\]
则 \(\lambda\in L\)，因为 \(L\) 对加法和减法封闭。于是：

\[
a' = a + \lambda,
\qquad \lambda\in L.
\]

因此必然有

\[
a' - a \in L,
\]
也就是

\[
a' \equiv a \pmod L.
\]

\begin{proposition}
在 Gentry09 的上述抽象写法中，总有
\[
a' \equiv a \pmod L.
\]
\end{proposition}

\begin{proof}
由
\[
a' = a + (\ell-\ell')
\]
且 \(\ell,\ell'\in L\)，可知 \(\ell-\ell'\in L\)。因此
\[
a'-a \in L,
\]
即
\[
a' \equiv a \pmod L.
\]
\end{proof}

\subsection{这一步的意义}

这一步告诉我们：解密输出 \(a'\) 与原始 \(a\) 一定属于同一个模格陪集。

但这还不是最终正确性，因为消息恢复不是模格，而是模 \(2\)。

所以现在的关键问题变成：

\begin{quote}
仅有
\[
a'-a\in L
\]
是否足以推出
\[
a' \equiv a \pmod 2 \ ?
\]
\end{quote}

答案是：\textbf{一般不足以推出。}

\section{为什么“\(a'-a\in L\)”一般不够}

\subsection{模 \(2\) 正确性真正需要什么}

由
\[
a' = a+\lambda
\]
可得
\[
a' \bmod 2 = (a+\lambda)\bmod 2 = (a\bmod 2)+(\lambda\bmod 2).
\]

因此，要推出
\[
a' \bmod 2 = a \bmod 2,
\]
必须证明
\[
\lambda \equiv 0 \pmod 2.
\]

换言之，仅仅知道 \(\lambda\in L\) 还不够，还必须进一步知道这个格向量在模 \(2\) 下消失。

\subsection{反例}

\begin{example}
取一维格
\[
L=\mathbb Z.
\]
令
\[
a=0,\qquad \lambda=1\in L.
\]
则
\[
a' = a+\lambda = 1.
\]
于是
\[
a\bmod 2 = 0,\qquad a'\bmod 2 = 1.
\]

因此，虽然
\[
a'-a = 1 \in L,
\]
但消息却发生了变化。
\end{example}

这个反例说明：

\[
\boxed{
a'-a\in L \quad \text{一般不能推出} \quad a'\equiv a\pmod 2.
}
\]

所以如果有人只证明到“\(a'\) 与 \(a\) 模格同余”，那对 bit 解密正确性来说通常是不够的。

\section{两种足够的正确性路线}

为了保证最终恢复消息，需要更强条件。常见有两条路线。

\subsection{路线一：强形式，证明 \(a'=a\)}

这是最直观的路线。

若能证明
\[
\ell'=\ell,
\]
则由
\[
a' = c-\ell',\qquad c=\ell+a
\]
立刻得到
\[
a' = (\ell+a)-\ell = a.
\]

于是
\[
a'\bmod 2 = a\bmod 2 = (m+2e)\bmod 2 = m.
\]

\begin{proposition}
若解密时恢复出的格点满足
\[
\ell'=\ell,
\]
则必有
\[
a'=a,
\]
因此消息正确恢复。
\end{proposition}

\begin{proof}
由
\[
a' = c-\ell'
\]
和
\[
c=\ell+a
\]
代入得
\[
a' = (\ell+a)-\ell' = a+(\ell-\ell').
\]
若 \(\ell'=\ell\)，则
\[
a'=a.
\]
于是
\[
a'\bmod 2 = a\bmod 2 = (m+2e)\bmod 2 = m.
\]
\end{proof}

\subsection{这条路线的几何意义}

这条路线要求私钥优质基足够好，使得解密算法恢复出的格点 \(\ell'\) 正好是加密分解中对应的那个格点 \(\ell\)。这正是 Babai 算法与最近格点问题进入正确性分析的地方。

若 \(a\) 很小，则
\[
c=\ell+a
\]
表示 \(c\) 离格点 \(\ell\) 很近。若 \(\ell\) 是唯一最近格点，且优质基保证恢复成功，就会得到 \(\ell'=\ell\)，从而 \(a'=a\)。

这是一条\textbf{很强、很干净}的正确性路径。

\subsection{路线二：弱形式，只需证明模 \(2\) 不变}

对 bit 解密而言，其实不必要求 \(a'=a\)。只要能证明
\[
a' \equiv a \pmod 2,
\]
就已经足够。

而这又可由更弱的条件推出，例如
\[
a'-a \in 2\mathbb Z^m,
\]
或更一般地，
\[
a'-a
\]
落在某个模 \(2\) 后消失的理想/子格中。

\begin{proposition}
若
\[
a'-a \in 2\mathbb Z^m,
\]
则
\[
a' \equiv a \pmod 2,
\]
从而消息正确恢复。
\end{proposition}

\begin{proof}
由
\[
a'-a \in 2\mathbb Z^m
\]
可知 \(a'-a\) 的每个坐标都是偶数，因此
\[
a'-a \equiv 0 \pmod 2.
\]
于是
\[
a' \equiv a \pmod 2.
\]
又因为
\[
a = m+2e,
\]
所以
\[
a \equiv m \pmod 2.
\]
综上，
\[
a' \equiv m \pmod 2,
\]
消息正确恢复。
\end{proof}

\subsection{这条路线比前一条弱，但已经够用}

需要特别强调：

\begin{itemize}[leftmargin=2em]
    \item “证明 \(a'=a\)”是更强结论；
    \item “证明 \(a'\equiv a\pmod 2\)”是恢复消息所需的最小层面结论之一；
    \item 因此，若方案的结构保证
    \[
    a'-a
    \]
    自动落在模 \(2\) 消失的理想或子格中，那么完全没必要再追求 \(a'=a\)。
\end{itemize}

从密码正确性的角度看，弱路线已经足够。

\section{如何理解这两条路线的关系}

\subsection{强路线推出弱路线}

显然，若
\[
a'=a,
\]
则自动有
\[
a'\equiv a \pmod 2.
\]

因此：

\[
\text{强路线} \Longrightarrow \text{弱路线}.
\]

\subsection{弱路线不推出强路线}

反过来不成立。完全可能有
\[
a' \neq a,
\]
但仍有
\[
a' \equiv a \pmod 2.
\]

例如一维情形下，
\[
a=1,\qquad a'=3.
\]
则
\[
a'\neq a,
\]
但
\[
a' \equiv a \pmod 2.
\]

因此，消息正确恢复并不要求逐点恢复原始 \(a\)。

\subsection{复习时最容易混淆的地方}

很多讲解一上来就说“解密恢复 \(a\)”；这种说法在直觉层面有帮助，但若不加说明，容易让人误以为这是唯一必要的证明目标。

实际上更准确的区分应该是：

\begin{center}
\begin{tabular}{lll}
\toprule
层次 & 目标 & 是否必需 \\
\midrule
强层次 & 证明 \(a'=a\) & 非必需，但很强很清楚 \\
弱层次 & 证明 \(a'\equiv a\pmod 2\) & 对 bit 正确性已足够 \\
更弱层次 & 仅证明 \(a'-a\in L\) & 一般不够 \\
\bottomrule
\end{tabular}
\end{center}

\section{推荐的严谨表述}

如果要在复习笔记、讲义或考试答案中用一句话准确描述 \(a\) 与 \(a'\) 的关系，推荐使用下面这版：

\begin{quote}
由加密与解密定义可知，解密输出 \(a'\) 与原始 \(a\) 总满足
\[
a'-a\in L.
\]
但这一般不足以推出消息正确。若进一步证明解密恢复出与加密对应的同一格点，则有 \(a'=a\)；更一般地，只要证明 \(a'-a\) 落在模 \(2\) 消失的理想或偶子格中，就可推出
\[
a' \equiv a \pmod 2,
\]
从而正确恢复 bit 消息。
\end{quote}

这句话的好处在于：

\begin{itemize}[leftmargin=2em]
    \item 先写出总是成立的基础关系；
    \item 再指出这还不够；
    \item 最后明确给出两种足够条件。
\end{itemize}

\section{最终总结}

本文专门整理了 Gentry09 中 \(a\) 与 \(a'\) 的关系。核心结论如下。

\subsection*{结论 1}
总有
\[
a' = a + (\ell-\ell'),
\]
因此
\[
a'-a \in L,
\qquad\text{即}\qquad
a' \equiv a \pmod L.
\]

\subsection*{结论 2}
仅有
\[
a'-a\in L
\]
一般不足以推出 bit 解密正确，因为一般格向量不一定在模 \(2\) 下消失。

\subsection*{结论 3}
若能证明
\[
\ell'=\ell,
\]
则得到最强结论
\[
a'=a,
\]
从而消息正确恢复。

\subsection*{结论 4}
对 bit 正确性而言，其实不必要求 \(a'=a\)。只要能证明
\[
a'-a \equiv 0 \pmod 2,
\]
例如
\[
a'-a\in 2\mathbb Z^m
\]
或落在模 \(2\) 后消失的理想/子格中，就已经足够。

\subsection*{一句话版}
\begin{quote}
\(a'\) 与 \(a\) 总是模格同余，但要恢复消息，还必须进一步证明它们在模 \(2\) 意义下相同；证明 \(a'=a\) 是一种更强但非唯一必要的路线。
\end{quote}

\end{document}